\documentclass[../main.tex]{subfiles}
\begin{document}
%==========================================TIÊU ĐỀ TẬP========================================
\begin{table}[h]
  \begin{adjustwidth}{-1cm}{}
    \begin{tabular}{>{\centering\arraybackslash}p{6cm}|>{\raggedright\arraybackslash}p{12.5cm}}
      \multirow{5}{6cm}
      % Cột bên trái: ÉPISODE và số tập
      {\\[-40pt]
        \flaregothic\fontsize{20pt}{20pt}\selectfont \centering
        \color{eptype}ÉPISODE\color{black}\\
        \burbank\fontsize{72pt}{72pt}\selectfont
        \color{epnum}27\color{black}\\[6pt]
        \cafeteria\fontsize{20pt}{20pt}\selectfont\color{epnum}(2-10)\color{black}
        \\[18pt]
        \flaregothic\fontsize{18pt}{18pt}\selectfont
        \color{epattr}ÉPISODE PERDU
        \color{black}
      }
       &
      % Dòng thứ nhất: tên tập tiếng Nhật
      {\fotskip\fontsize{18pt}{18pt}\selectfont 〇〇\ruby{飼}{か}った\ruby{結}{けっ}\ruby{果}{か}wwwwwwwww
      } \\[6pt]
      % Dòng thứ hai: tên tập tiếng Anh
       & {\cafeteria\fontsize{18pt}{18pt}\selectfont Keeping a 〇〇!!}                           \\[4pt]
      % Dòng thứ ba: tên tập tiếng Việt
       & {\vnmsans\fontsize{18pt}{18pt}\selectfont Thú cưng Frappuccino}                       \\[8pt]
      % Dòng thứ tư: Nhân vật xuất hiện
       & {\caviar{\fontsize{14pt}{14pt}\selectfont Nhân vật: \emp{kuon1} \emp{ryuuhou1} \emp{achan} \emp{roboco} \emp{miko} \emp{fubuki} \emp{matsuri} \emp{mio}}} \\[8pt]
      % Dòng cuối cùng: Thời gian diễn ra của tập (thường sẽ là ngày phát hành)
       & {\continuum\fontsize{16pt}{16pt}\selectfont Ngày phát hành: 10/11/2019}               \\[18pt]
    \end{tabular}
  \end{adjustwidth}
\end{table}
%=========================================TRUYỆN CHÍNH========================================
\color{black}

\txtbx[2]{{\color{vol} \uline{Tóm tắt:}} Nàng Cổ động Matsuri nằng nặc đòi nuôi một ``con'' Frappuccino khổng lồ trong văn phòng. Trong khi Kuon và Fubuki ban đầu phản đối nhưng sau đó lại bị cuốn vào vòng xoáy tưởng tượng điên rồ, nàng Sói Mio lại là người duy nhất cố gắng giữ tỉnh táo để rồi nhận cái kết đắng lòng.}

\pov{Ookami Mio}

Đó là một buổi sáng vốn dĩ rất yên tĩnh trong văn phòng \hll. Mọi người đều đang bận rộn với công việc của mình, còn tôi thì đang tận hưởng chút thời gian rảnh rỗi hiếm hoi. Tôi ngồi vắt vẻo trên ghế, một tay lướt điện thoại, một tay nhâm nhi miếng bánh gạo giòn rụm mà anh Kuon vừa mang đến.

\img{2-11a}

Thế nhưng, bầu không khí êm đềm ấy chẳng kéo dài được bao lâu. Một giọng nói khẩn thiết, cao vút của Matsuri bất ngờ vang lên, phá tan sự tĩnh lặng: ``Đi mà! Cho tôi nuôi nó đi mà\~{}!''

Tôi tò mò ngước nhìn về phía đó. Chị ấy đang đứng nài nỉ Fubuki một cách không biết mệt mỏi. Cô bạn Cáo trắng thì khoanh tay, gương mặt nghiêm nghị hiếm thấy như thể đang cố gắng giữ vững lập trường cuối cùng của mình. Đứng tựa vào bàn gần đó là anh Kuon, gương mặt anh ấy hiện rõ vẻ chán nản và bất lực tột độ.

Ngồi cạnh tôi, Ryuuhou cũng ngừng dán mắt vào máy game cầm tay, con bé chống cằm nhìn hội tiền bối rồi thầm thì: ``Này chị Mio, em cá là chị Matsuri lại sắp bày ra trò gì đó ảo ma lắm cho xem.''

Tôi gật đầu đồng tình, vừa nhai bánh vừa tự hỏi: ``Nuôi thú cưng trong văn phòng này sao? Liệu có ổn không đây?''

Anh Tổng thở dài, lên tiếng đáp lại bằng giọng điệu không thể mệt mỏi hơn: ``Anh đã nói rồi, công ty không cấm mang thú cưng sang chơi, nhưng để nuôi hẳn một con tại văn phòng thì thực sự là nhiệm vụ bất khả thi đấy.''

Nhưng Matsuri-senpai nào có chịu bỏ cuộc. Chị ấy dậm chân, khăng khăng đòi bằng được: ``Xin hai người đấy! Chị thực sự muốn có một con cho riêng mình!''

Thấy chị ấy kiên quyết quá, anh Tổng đành phải xuống nước, tặc lưỡi hỏi: ``Thôi được rồi, cứ cho là chúng ta được phép nuôi đi. Thế rốt cuộc em muốn nuôi con gì?''

Chỉ chờ có thế, chị hít một hơi thật sâu rồi tuôn ra một tràng dài dằng dặc với tốc độ ánh sáng: ``Chị muốn nuôi một con frappuccino cỡ venti, loại extra-shot, kèm theo hạt dẻ, vani, hạnh nhân, caramel, thêm kem đánh, sốt caramel, sốt mocha, rắc thêm bột cà phê và thật nhiều kem sô-cô-la lên trên nữa!''

\begin{figure}[H]
  \centering
  \animategraphics[autoplay,loop,width=12.5cm]{30}{../../../../Image/Book2/2-11b/2-11b.}{1}{190}
  \caption{Tốc độ ``bắn'' chữ của Matsuri thực sự khiến người nghe phải chóng mặt.}
\end{figure}

Tôi đứng hình mất năm giây, miếng bánh gạo trong miệng suýt thì rơi ra ngoài. Ryuuhou thì mắt chữ O mồm chữ A: ``Cái gì cơ? Đó là thực đơn Starbucks chứ có phải thú cưng đâu!''

Kuon-senpai vẫn giữ được vẻ mặt bình thản đến đáng sợ. Anh ấy nhướng mày, mỉa mai: ``Tên dài ngoằng thế? Mà em đang mô tả một cốc cà phê ngon lành hơn là một sinh vật sống đấy.''

Fubuki cũng lắc đầu lia lịa: ``Không được đâu! Chúng ta không thể nuôi một con frappuccino cỡ venti\dots\ ờm\dots\ tên nó là gì tiếp theo ấy nhỉ?''

Như một chiếc máy ghi âm siêu hạng, anh Tổng lặp lại y hệt cái tên đó với tốc độ tương đương Matsuri-senpai khiến cả đám sững sờ: ``Extra-shot, hạt dẻ, vani, hạnh nhân, caramel, thêm kem đánh, sốt caramel, sốt mocha, rắc cà phê và kem sô-cô-la. Anh hiểu mà, tên dài quá nên Fubu-chan không nhớ nổi là phải.''

Tôi bật cười trước sự a dua của bà bạn mình: ``Bà còn chẳng nói ra nổi cái tên của `con' thú cưng đó mà cũng đòi nuôi sao!''

\ct

Tưởng rằng sau khi bị bác bỏ, câu chuyện sẽ kết thúc. Nhưng không, bộ ba đó bắt đầu bàn luận về ``con'' frappuccino như thể nó là một sinh vật sống thật sự. Tôi ngồi im thin thít, cố gắng tiêu hóa cái logic kỳ quái đang diễn ra trước mắt.

\begin{figure}[H]
  \centering
  \begin{subfigure}[b]{0.45\textwidth}
    \centering
    \includegraphics[width=8cm]{../../../../Image/Book2/2-11c1.png}
  \end{subfigure}
  \quad
  \begin{subfigure}[b]{0.45\textwidth}
    \centering
    \includegraphics[width=8cm]{../../../../Image/Book2/2-11c2.png}
  \end{subfigure}
\end{figure}

Fubuki hỏi với vẻ mặt cực kỳ nghiêm túc: ``Thế bà có chắc là mình sẽ chăm sóc nó mỗi ngày không?'' 
Tôi thầm nghĩ: ``\textit{Xạo ke vừa thôi, cà phê mà cũng cần chăm sóc sao!}''

Kuon-senpai còn bồi thêm một câu xanh rờn: ``Rồi còn chuyện cho nó ăn, cho nó uống, rồi cả việc dọn `sản vật' của nó nữa\dots\ Em có làm nổi không?''
Tôi suýt ngã ngửa: ``\textit{Cà phê chứ có phải chó mèo đâu mà có `sản vật' để dọn!}''

Ryuuhou thì cười nắc nẻ: ``Anh nhìn kìa chị Mio, Onii-san đang nhập vai quá mức rồi!''

Matsuri lúc này hăng máu hẳn lên, chị ấy dõng dạc tuyên bố dù giọng hơi run: ``Tôi sẽ\dots\ ờ\dots\ đem nó đi triển lãm để tạo ra một điểm thu hút du lịch mới cho công ty!''

Tôi gật gù trong lòng: ``Hóa ra là vì tiền thôi sao!''

Nhưng lập tức dội gáo nước lạnh: ``\textit{Lý do này không ổn. Ai mà thèm bỏ tiền ra xem một cốc cà phê trong tủ kính khi họ có thể tự mua một cốc y hệt ở cửa hàng bên cạnh chứ Kuon-senpai?}''

Chị Cổ động vẫn không chịu thua: ``Tôi sẽ sáng tạo ra nhiều cách để mọi người thấy được vẻ đẹp của nó trong môi trường hoang dã!''

Tôi cạn lời luôn: ``Frappuccino mà cũng có môi trường hoang dã sao!?''

\img{2-11d}

Chị Cổ động bắt đầu vẽ ra một khung cảnh tưởng tượng. Trong đầu chị ấy, ``con'' frappuccino đang đứng hiên ngang trên một màng mỏng, còn tôi, Fubuki và Kuon-senpai thì đang ngước nhìn lên với vẻ đầy sùng bái.

``Đây là vuốt của nó nè!'' Matsuri-senpai hăng hái chỉ vào cái cốc ảo ảnh.
Tôi lẩm bẩm: ``Vuốt cái hú gì cơ!?''

Fubuki gật gù tán thưởng: ``Ừm, nhìn kỹ thì cũng giống thật!''

Kuon-senpai thì nhận xét chuyên môn: ``Anh thấy thiếu phần chân sau\dots\ nhưng nhìn chung ý tưởng này cũng khá triển vọng đấy.''

Tôi đưa tay ôm trán, cảm giác mình sắp phát điên đến nơi rồi. Hai người kia làm ơn tỉnh lại dùm tôi cái! Đừng có tiếp tay cho cái sự điên rồ này nữa!

\img{2-11e1}

Fubuki bỗng nảy ra ý tưởng mới: ``Hay là cho nó sống dưới đại dương đi?'' Tưởng tượng một cốc cà phê chìm giữa rặng san hô\dots\ Thôi xong, mấy người này hết thuốc chữa thật rồi.

Matsuri reo lên: ``Bà đúng là thiên tài mà Fubuki-chan!!''

Bất thình lình, cả ba người cùng quay phắt lại nhìn tôi với nụ cười nham hiểm: ``\dots Và chúng ta sẽ cần một người vào vai mỹ nhân ngư để bơi cùng nó. Mio chắc chắn là người phù hợp nhất rồi!''

Tôi giật bắn mình như bị điện giật, miếng bánh gạo trên tay rơi bộp xuống đất. Trong đầu tôi gào thét: ``Không đời nào! Chắc chắn họ đang cố tình lôi mình vào để làm trò cười đây mà! Phải nhẫn nhịn, không được để ý đến họ!''

\img{2-11f}

Nhưng Matsuri vẫn chưa tha cho tôi: ``Chị cũng muốn thấy Mio-chan bơi lội cùng cốc cà phê nữa!''

\img{2-11e2}

Fubuki thêm mắm dặm muối: ``Bộ đồ bơi của Mio hẳn là sẽ đáng yêu lắm!''

Kuon-senpai thì hào hứng: ``Xinh thế này thì phải chụp ảnh cháy máy mới thôi!''

Tôi quay ngoắt đi, úp mặt xuống ghế sô-pha, tay bịt chặt tai để không phải nghe thêm bất kỳ lời nhảm nhí nào nữa. Cốc cà phê, đồ bơi, đại dương\dots\ lũ người này điên hết thật rồi!

Nhưng ngay khi tôi tưởng mình sắp thoát khỏi cơn ác mộng này, cánh cửa phòng bật mở. Roboco-senpai bước vào với giọng nói trong trẻo: ``Này mọi người, nhìn xem chị có gì nè!''

Tôi hé mắt nhìn, thấy chị ấy tay cầm một chai nước rỗng, tay kia cầm một cốc cà phê. Roboco-senpai cười tươi: ``Một chai PET có con `pet' ở trên này nè!''

\img{2-11g}

Vừa nhìn thấy cốc cà phê, bản năng tôi bỗng trỗi dậy mạnh mẽ đến mức không thể kiểm soát nổi. Tôi bật dậy như một chiếc lò xo, đôi mắt trợn tròn, và tiếng hú vang dội cả văn phòng: ``\textbf{FRAPPUCCINO NHÉ!!!}''

Khoảnh khắc đó, tôi nhận ra mình đã thua cuộc một cách thê thảm. Tôi vừa muốn khóc vừa muốn cười vì sự yếu lòng của chính mình.

\img{2-11h}

Cả văn phòng nổ tung trong tiếng cười. Tôi ngồi thẫn thờ nhìn họ, trong lòng tự nhủ: ``Chuyện này thậm chí còn chẳng buồn cười tẹo nào cả\dots''

%==========================================TRUYỆN PHỤ=========================================
\omk
\pov{Hikawa Kuon}

Ngay khi nàng Sói đen nhận ra mình đã sập bẫy, Mio-chan đổ gục xuống sàn, hai tay đấm thình thịch xuống đất trong cơn uất ức tột cùng. Gương mặt em ấy đỏ gay, nước mắt ngắn nước dài nhìn cứ như sắp nổ tung đến nơi.

\img{2-11i}

Nhưng thay vì an ủi, cả lũ chúng tôi - bao gồm cả tôi, Matsuri, Roboco và Fubuki - đồng loạt nhảy một điệu ăn mừng vô cùng lố bịch xung quanh em ấy. Thậm chí Ryuuhou còn lấy điện thoại ra quay phim lại và cười sằng sặc: ``Anh nhìn kìa Onii-san, biểu cảm của chị Mio đúng là cực phẩm luôn!''

Tôi vừa nhảy vừa hú hét, thực sự cảm thấy công sức bày mưu tính kế từ sáng đến giờ thật xứng đáng.

Sau khi cười đùa chán chê, Fubuki mới tỏ ra hối lỗi một chút. Cô nàng tiến lại gần, đưa cốc frappuccino thật cho Mio-chan: ``Thôi nào Mio-chan, cốc này là của bà đó. Xem như quà bù đắp cho công sức bà đã kiên trì suốt cả buổi sáng!''

Mio-chan hậm hực nhận lấy cốc nước, cơn giận cũng nguôi đi phần nào khi hơi lạnh của cà phê lan tỏa.

Nhưng cuộc vui vẫn chưa kết thúc ở đó. Tôi vỗ tay thu hút sự chú ý: ``Nào mọi người! Đã nói là có thưởng thì phải thực hiện. Tối nay anh mời tất cả đi ăn một bữa ra trò với đám bạn của anh!''

Fubuki và Matsuri reo lên đầy phấn khích. Roboco cũng hăng hái: ``Được Kuon-senpai mời thì phải ăn cho sập tiệm mới thôi!''

Tôi liếc nhìn nàng Người máy, châm chọc: ``Thế Robo-PON ăn được gì? Đinh tán hay dầu máy?'' Roboco liền phồng má: ``Em ăn được đồ người bình thường! Và đã bảo là em không có PON mà!''

Ryuuhou đứng bên cạnh kéo áo tôi: ``Onii-san, em cũng đi nữa! Em phải trông chừng không cho anh làm gì xấu với các chị ấy!''

Tôi cốc đầu con bé một cái: ``Nghịch vừa thôi, cho em đi để em ăn hết phần của anh à?''

Nhìn sang Mio-chan, em ấy vẫn đang ngồi thu lu một góc với vẻ mặt tủi thân vì nghĩ mình không được mời. Tôi khẽ mỉm cười, tiến lại gần xoa đầu em ấy: ``Này Mio-chan, đi năm người thì vui hơn bốn, mà đi sáu người thì lại càng vui hơn. Em đi cùng bọn anh luôn?''

Mio-chan ngẩng phắt đầu lên, đôi mắt long lanh hạnh phúc. Em ấy mừng đến phát khóc - lần này là khóc thật luôn - và liên tục cúi đầu cảm ơn tôi: ``Cảm ơn anh Kuon-senpai! Anh đúng là người công bằng nhất thế gian luôn!''

Thế là cả hội kéo nhau đi ăn tối. Dù ai thắng ai thua, thì rốt cuộc tất cả đều kết thúc bằng những tiếng cười rộn rã. Đúng như tôi đã nói:

``Đông người thì lúc nào cũng vui hơn mà, đúng không nào?''

\end{document}